\section{Experimental Setup}

\subsection{Models}

We evaluate twelve models spanning a wide range of sizes, providers, and access modes on the full 406-item benchmark (Table~\ref{tab:models}, Appendix~\ref{app:models}): Gemini 2.5 Flash and Gemini 2.5 Pro (Google APIs); Qwen3-32B, LLaMA-3.3-70B, Llama 4 Scout 17B, Kimi K2, GPT-OSS 120B, GPT-OSS 20B, and DeepSeek-R1-Distill-Llama-70B (Groq API); and LLaMA-3-8B, Mistral-7B, and Gemma 4 E4B run locally via Ollama v0.6.5 in 4-bit quantization (Q4\_K\_M) on an RTX 4060 (8~GB). All models were evaluated under identical zero-shot conditions with no fine-tuning or prompt adaptation; exact checkpoint identifiers are documented in the evaluation code.

\subsection{Prompting Strategy}

Our evaluation follows the extractive QA paradigm established by SQuAD \citep{rajpurkar2016squad}, in the same spirit as DROP's discrete-reasoning-over-text framing \citep{dua2019drop}, with context passages provided directly to the model under zero-shot, context-only constraints. This design isolates the model's ability to reason about regulatory text rather than recall memorized facts, making the benchmark robust to training data contamination. All models received a system prompt establishing the context-only constraint: ``You are an expert in Indian financial regulation and policy. Answer questions using ONLY the provided context passage. Do not use any external knowledge. Be concise and precise. Give only the answer---no preamble.'' Task-specific user prompts provided appropriate formatting instructions for each task type. For contradiction detection, both passages were labelled explicitly as `Passage A / Passage B'. For numerical reasoning, models were instructed to show calculation steps and include units---an instruction in tension with the system prompt's ``no preamble'' constraint, which contributes to the verbosity-related scoring penalty quantified in Appendix~\ref{app:judge}. All models were evaluated under identical prompting and decoding settings (temperature = 0.0; completion budget 200 tokens for API models, 300 for local models).

We evaluate exclusively under zero-shot conditions, as this most closely reflects practical deployment where users query models without domain-specific priming, and because it eliminates confounds from example selection strategy and ordering effects. Few-shot results for four top models are reported in Appendix~\ref{app:fewshot}.

\subsection{Scoring}
\label{sec:scoring}

Answers were scored using a multi-stage matching procedure applied in sequence: (1) exact match after case-normalization and punctuation stripping; (2) fuzzy token match using RapidFuzz \texttt{token\_set\_ratio} $\geq$ 0.72, applied when exact match fails, with the 0.72 threshold established by manual inspection of 20 borderline cases and validated against adjacent thresholds (0.65 and 0.80); (3) numerical extraction match, where the item is scored correct if the set of numbers extracted from both the reference and prediction are identical (handling currency symbols, comma separators, and units); and (4) Yes/No match for contradiction detection, comparing the leading word of the prediction to the reference.

To measure the false-negative rate of the automated pipeline, we applied a secondary LLM-as-judge validation using Gemini 2.5 Flash over items marked incorrect on NUM, REG, and TMP tasks across the twelve models. Of 935 flagged predictions, 874 were judged; the remaining 61 (58 empty responses and 3 API failures) are unambiguously incorrect and required no semantic review. CON items use exact Yes/No matching and were not judged. Judge reliability was confirmed on a 28-item stratified manual audit (89.3\% accuracy). Full results are in Appendix~\ref{app:judge}; primary Table~\ref{tab:mainresults} scores use the strict automated pipeline. DeepSeek-R1-Distill-Llama-70B was accessed through an endpoint that returns final answers without exposed reasoning traces; its stored predictions contain no chain-of-thought markup, and its internal deliberation shares the fixed completion budget---a constraint we return to in Section~\ref{sec:discussion}.

\subsection{Human Reference Point}
\label{sec:humanbaseline}

To anchor model scores against human performance, a single non-specialist evaluator (a university-educated adult with no professional background in finance or law) independently answered 60 benchmark items sampled from the medium (40) and hard (20) difficulty strata across all four task types (REG 11, NUM 16, CON 17, TMP 16). The evaluator worked from the same context passages and questions given to the models, without access to reference answers or external resources, and responses were scored by the identical automated pipeline of Section~\ref{sec:scoring}. The evaluator scored 80.0\% (48/60; 95\% Wilson CI [68.2\%, 88.2\%]): REG 100\% (11/11), TMP 87.5\% (14/16), CON 82.4\% (14/17), NUM 56.2\% (9/16). Because the sample deliberately excludes easy items, this is a reference point for the benchmark's harder portion, not an estimate of whole-benchmark human accuracy; with one evaluator it reflects a single careful non-specialist rather than a population. Section~\ref{sec:humananalysis} compares each model to this reference on the identical 60 items using paired bootstrap tests.

